\documentclass[10pt, a4paper]{article}

% Packages for layout and styling
\usepackage[utf8]{inputenc}
\usepackage[margin=0.75in]{geometry}
\usepackage{enumitem}
\usepackage{hyperref}
\usepackage{titlesec}
\usepackage{parskip}
\usepackage{etoolbox}

\input{personal_details.tex}

\ifdefined\showPrivate
\else
  \newcommand{\showPrivate}{true}
\fi

\newbool{showInfo}
\ifdefstring{\showPrivate}{true}{\booltrue{showInfo}}{\boolfalse{showInfo}}
\newcommand{\personal}[1]{\ifbool{showInfo}{#1}{\rule{2cm}{0.5pt}}}


\titleformat{\section}{\large\bfseries}{}{0em}{}[\titlerule]
\titlespacing{\section}{0pt}{10pt}{5pt}

\begin{document}

\begin{center}
    {\LARGE \textbf{Oliver Zimmermann}} \\
    \vspace{2pt}
    \street, \city \(\mid\) \phone \(\mid\) \email \\
    Date of birth: \dateofbirth \(\mid\) Place of birth: \placeofbirth \(\mid\) Nationality: \nationality
\end{center}

\section*{Education}

\textbf{ETH Zürich} \hfill Zurich, Switzerland \\
Master of Science: Data Science \hfill since September 2025
\begin{itemize}[noitemsep, topsep=0pt]
    \item \textbf{Specialization} Mathematical foundations of ML; AI for Science (e.g. Causal Inference, Data Science for Genomics, Computational Chemistry, AI in Science and Engineering)
\end{itemize}

\textbf{Hasso-Plattner-Institut - Universität Potsdam} \hfill Potsdam, Germany \\
Bachelor of Science: IT-Systems Engineering \hfill September 2021 - August 2025
\begin{itemize}[noitemsep, topsep=0pt]
    \item \textbf{Grade:} 1.2 with distinction (German grading system - 1.0 is best)
    \item \textbf{Thesis:} Decentralized Multi-Agent Reinforcement Learning for Cooperative Navigation with
Information Exchange in Simulated Ad Hoc Networks (Grade 1.0)
    \item \textbf{Specialization:} Machine learning and statistics (statistics and probability theory, introduction to
probabilistic machine learning, medical machine learning seminar, reinforcement learning project)
\end{itemize}

\textbf{Technische Universität Berlin} \hfill Berlin, Germany \\
Bachelor of Science: Natural Sciences in the Information Society \hfill October 2024 - August 2025
\begin{itemize}[noitemsep, topsep=0pt]
    \item Fundamental courses in mathematics accounting for 46 additional ECTS (1/4 of a full degree)
\end{itemize}

\textbf{Georg-Cantor-Gymnasium (STEM-specialized high school)} \hfill Halle (Saale), Germany \\
German Abitur: 1.0 (best possible grade) \hfill September 2013 -- July 2021
\begin{itemize}[noitemsep, topsep=0pt]
    \item Advanced courses in Physics, Mathematics, and Computer Science
    \item Recipient of the top school award for academic excellence and social involvement (granted to only one
graduate annually)
    \item Winner of multiple national and international STEM competitions (Mathematics, Physics, Chemistry, Engineering)
\end{itemize}

\section*{Experience}

\textbf{QuantCo} \hfill Berlin, Germany \\
Data and Software Engineer (working student) \hfill February 2025 -- August 2025
\begin{itemize}[noitemsep, topsep=0pt]
    \item Improved data processing pipeline by implementing streamed data analytics and efficient data archiving (Kafka, Parquet, DuckDB)
    \item Assisted in deploying core product as a cloud solution (Kubernetes, Terraform, Helm, AWS)
    \item Simplified multi-repository demos by providing virtualized environments and reusable code templates (DockerCompose, Python)
\end{itemize}

\textbf{Hasso-Plattner-Institut} \hfill Potsdam, Germany \\
Student Assistant (AI and Sustainability group) \hfill October 2023 -- August 2025
\begin{itemize}[noitemsep, topsep=2pt]
    \item Assisted in preparing the "Statistics" and "Introduction to Probabilistic Machine Learning" courses
    \item Helped design a Massive Open Online Course (MOOC) on probabilistic machine learning, which has
been taken by over 3,000 learners worldwide \footnote{Course website: \url{https://open.hpi.de/courses/probabilisticai2023}}
\end{itemize}

Student Assistant (AI and Intelligent Systems group) \hfill April 2023 -- September 2023
\begin{itemize}[noitemsep, topsep=2pt]
    \item Contributed to a reusable NLP template for high-performance compute clusters used by master's and
PhD students in research
    \item Benchmarked various NLP models based on training configurations (e.g., numeric precision) and
created a report with best practices for efficient model fine-tuning
\end{itemize}

\textbf{Hospital Ernst von Bergmann} \hfill Potsdam, Germany \\
Internship IT \hfill April 2023 -- September 2023
\begin{itemize}[noitemsep, topsep=2pt]
    \item Collected, cleaned, and analyzed clinical data to improve hospital processes and predict capacity needs
    \item Applied process mining techniques to identify bottlenecks in hospital procedures
\end{itemize}

\section*{Scholarships \& Awards (selection)}
    \textbf{European Space Agency - CanSat Competition, Grand Prize Winner} \hfill June 2019 \\
    I represented Germany as part of the national team at the 2019 European CanSat Competition, hosted by
the European Space Agency in Italy. We developed a miniature satellite that performed a stabilized,
autonomous landing and transformed into a quadcopter for additional data collection, acting as an
exploratory space probe. Our team won the European Championship, competing against 20 international
teams. Our technical report is featured on the ESA website.

\textbf{Sholarship: German Academic Scholarship Foundation} \hfill since December 2022 \\
    I receive a scholarship from the German Academic Scholarship Foundation for exceptional academic
achievements and a positive contribution to society. It is Germany's most prestigious scholarship, awarded to less then 0.5\% of students.

\textbf{Distinguished STEM Certificate} \hfill July 2021 \\
    Awarded for exceptional performance in the national mathematics, physics, chemistry, and computer science
Olympiads in Germany and other science competitions (e.g. Jugend forscht).


\section*{Projects (selection)}
\textbf{AI for Science} \hfill January 2026 \\
My final project for the AI in Science and Engineering Course at ETH Zurich. I analyzed spectral bias in physics-informed neural networks (PINNs)
using techniques from loss landscape visualization. I also trained a Fourier Neural Operator (FNO) to solve a dynamical system and showed that
this architecture is indeed not resolution invariant, despite many claims in the current literature. Finally, I contributed to the codebase of the
Geometry-Aware Operator Transformer (GAOT), implementing an adaptive radii sampling strategy for unstructured mesh data. \footnote{The source code for the AI in Science project is available on GitHub \url{https://github.com/OliverZim/aise_project}}\\


\textbf{DEFIANCE} \hfill September 2023 -- August 2024 \\
    Together with a team of 7 students, I developed the DEFIANCE extension for the network simulator ns-3.
This extension enables research in online reinforcement learning within multi-agent environments and
includes realistically simulated wireless communication between agents. The project was open-sourced and
is now an official extension of the ns-3 simulator. It has been included in the project portfolio of the Open
6G Hub of the German Federal Ministry of Education and Research and I was invited to present it at the
Berlin 6G Conference 2024. This work has lead to a publication in ICNS3 in 2025 (doi: \href{https://doi.org/10.1145/3747204.3747205}{10.1145/3747204.3747205}).\\
\footnote{The source code for the DEFIANCE project is available on GitHub \url{https://github.com/DEFIANCE-project/ns3-DEFIANCE}}

\textbf{REMEMBER} \hfill April 2023 -- October 2023 \\
In this project I applied deep learning techniques for Alzheimer's disease classification using 3D MRI data
from the Alzheimer's Disease Neuroimaging Initiative and the UK Biobank. My experiments focused on
semi-supervised learning as this allowed me to utilize the large fraction of unlabeled medical data for the
classification task. I compared different model architectures and pretraining methods (e.g. BYOL,
SIMCLR) and achieved good results with an adaptation of the Masked Autoencoder architecture.

\section*{Interests \& Skills}
\textbf{Languages:} German (native), English (fluent, C1), Python, Julia, C++\\
\textbf{Professional Skills:} Public speaking, presenting, team coordination \\
\textbf{Personal Interests:} Swimming, athletics, ballroom dance (standard and Latin), DJing

\end{document}